\section{Notation}\label{notation-1}
\begingroup
\setlength{\parskip}{3pt}

\textbf{G}\textsubscript{\textbf{R}}\textbf{ } := GB of \emph{raw} capacity

\textbf{G}\textsubscript{\textbf{A}}\textbf{ } := GB of \emph{available} capacity

\textbf{G}\textsubscript{\textbf{U}} := GB of potentially usable capacity under the protection-only approximation

\textbf{C}\textsubscript{\textbf{R}} := Material cost of \emph{raw} capacity, expressed as \$:G\textsubscript{R}

\textbf{C}\textsubscript{\textbf{U}} := Cost with field uplift per actually usable GB

\textbf{D} = \{D\textsubscript{0}, D\textsubscript{1}, D\textsubscript{2}\} -- HDD types contained in Nodes

\begin{quote}
\textbf{D}\textsubscript{\textbf{0}} := 320GB HDD (Gen 4 HDD)

\textbf{D}\textsubscript{\textbf{1}} := 500GB HDD (Gen 4\textsuperscript{\ensuremath{\prime}} HDD)

\textbf{D}\textsubscript{\textbf{2}} := 500GB HDD (Gen 5 HDD)
\end{quote}

\textbf{N} = \{N\textsubscript{0}, N\textsubscript{1}, N\textsubscript{2}\} -- Node types, represented by tuples of their installed disks

\begin{quote}
\textbf{N}\textsubscript{\textbf{0}} = (D\textsubscript{0}, D\textsubscript{0}, D\textsubscript{0}, D\textsubscript{0}) (Gen 4 node)

\textbf{N}\textsubscript{\textbf{1}} = (D\textsubscript{1}, D\textsubscript{1}, D\textsubscript{1}, D\textsubscript{1}) (Gen 4\textsuperscript{\ensuremath{\prime}} node)

\textbf{N}\textsubscript{\textbf{2}} = (D\textsubscript{2}, D\textsubscript{2}, D\textsubscript{2}, D\textsubscript{2}, D\textsubscript{2}, D\textsubscript{2}) (Gen 5 node)
\end{quote}

\textbf{H} = \{H\textsubscript{0}, H\textsubscript{1}, H\textsubscript{2}, H\textsubscript{3}\} -- Fixed overhead, expressed in GB

\begin{quote}
\textbf{H}\textsubscript{\textbf{0}} := Regeneration buffer (``regen buffer'' in the CLI)

\textbf{H}\textsubscript{\textbf{1}} := Database overhead (``system buffer'' in the CLI)

\textbf{H}\textsubscript{\textbf{2}}\textbf{ }:= FilePool + Platform (Linux + Filesystem) overhead (``system resources'' in CLI)

\textbf{H}\textsubscript{\textbf{3}} := Audit and metadata overhead
\end{quote}

\textbf{L} := ``Cluster object'' count limit, anticipated to be 25M:HDD in Hapkin

\textbf{Z} := Number of HDD in the modeled cluster or homogeneous partition

\textbf{P} = \{P\textsubscript{cpm}, P\textsubscript{cpp}\} Data protection scheme, over the interval (0,1{]}, where P\textsubscript{cpm} := CPM and P\textsubscript{cpp} := CPP. Generally, \emph{P}\textsubscript{\emph{x}} = \emph{m} / (\emph{m} + \emph{n}), where there are \emph{m} data fragments and \emph{n} protection fragments.

\textbf{R} := Titan rack material cost + miscellaneous minor components (cables, etc.), in US\$

\textbf{W } := G4 Switch material cost (i.e. 2 switches) + 2 PDP + 2 PDU + cables, etc. per cube

\textbf{T} := Total Solution Cost, expressed in US\$

\textbf{A} := ASP per actually usable GB

\(\boldsymbol{j_i}\) := Number of physical nodes of type \(i\)

\textbf{u } := Fractional field uplift; for example, 10\% means \(u=0.10\)

A \textbf{user object} is the combination of a CDF plus associated blob(s) (typically a user file)

A \textbf{cluster object} is a single CDF or blob fragment

\textbf{Q}\textsubscript{\textbf{obj}} : = number of user objects in the system partition

\textbf{S}\textsubscript{\textbf{cdf}} := Average CDF size, which can be considered a constant of 4 KiB for our purposes

\textbf{S}\textsubscript{\textbf{blob}} := Average blob size (i.e. the size of the user object being ingested into the cluster)

\par\endgroup

\(C'_R\) denotes uplifted raw cost per GB; \(A_R\) denotes raw selling price per GB. \(T'=T(1+u)\) and \(T''=T'/(1-GM)\) are total uplifted cost and total selling price.

\(N=\sum_i j_i\) counts physical nodes. \(U=j_0+j_1+2j_2\) is the
occupied rack space in 1U units. \(C\) is the actual number of populated
cubes, with \(N_c\) physical nodes and \(G_{A,c}\) available GB in cube \(c\).
\(G'_U\) accounts for average blob and CDF sizes before object-count limits;
\(G_{U,Actual}\) also accounts for those limits. Both count logical blob
and CDF bytes once, excluding protection copies.

\subsection{Notes}\label{notes-3}

\begin{enumerate}
\def\labelenumi{\arabic{enumi}.}
\tightlist
\item
  All costs and prices are in US\$.
\item
  \label{note:normalization}A physical Gen 4 or Gen 4\textsuperscript{\ensuremath{\prime}}
  node occupies 1U; a physical Gen 5 node occupies 2U. Physical nodes are used
  for disk counts and protection groups. Rack-space units \(U\) are used
  for enclosure cost, so a 2U node is not counted twice when calculating raw capacity.
\item
  In earlier versions of the model we factored out the effect of CDF storage and object size to simplify the model. However, we subsequently decided that this simplification masked too much detail. As of v0.50 of the model, we have added consideration for these two factors.
\item
  Cost of minor elements, e.g. cables, are factored into component cost of switches, nodes and rack(s).
\item
  We assume new installs.
\item
  We do not consider the effect of single instance store (SIS).
\item
  CDFs are written in a 1+\emph{n} protection scheme; i.e. we have 1+\emph{n} copies of each CDF.
\item
  Capacity units must be consistent throughout a calculation. A GB is
  \(10^9\) bytes and a TB is \(10^{12}\) bytes; a GiB is \(2^{30}\) bytes
  and a TiB is \(2^{40}\) bytes, following the IEC notation in
  \hyperref[ref:2]{[2]}. Thus 500 GB is approximately 465.66 GiB.
  The numerical value is 6.87\% lower when expressed in GiB; the physical
  capacity is unchanged. Convert blob and CDF sizes into the same units as
  capacity before applying the object-count equations; for example,
  4 KiB is \(4096/10^9\) GB. The presentation in \hyperref[ref:8]{[8]}
  uses decimal units for comparison with advertised disk capacities.
\end{enumerate}
